\begin{center}
    {\LARGE \textbf{2021无机化学与物理化学}} \\[2ex]
    {\large 时量180分钟 \quad 满分150分}
\end{center}

\vspace{0.5cm}
\noindent\textbf{注意事项:}
\begin{enumerate}[label=\arabic*., parsep=0pt, itemsep=0pt]
    \item 答题前,考生先将自己的姓名、学号、考场号、座位号填写在答题卡上,将条形码准确粘贴在考生信息条形码粘贴区。
    \item 考试结束后,将本试卷与答题纸一并收回。
    \item 回答各个问题时,将答案写在答题纸上,写在本试卷上无效。
\end{enumerate}

\vspace{1cm}
\begin{center}
    {\Large \textbf{无机化学部分(共75分)}}
\end{center}

\vspace{0.5cm}
\noindent\textbf{一、选择题（每小题2分，共20分）}

\begin{enumerate}[label=\arabic*., itemsep=1.5ex]
    \item 下列卤素含氧酸中，酸性最强的是
    \begin{tasks}(4)
        \task \ce{HClO}
        \task \ce{HClO2}
        \task \ce{HClO3}
        \task \ce{HClO4}
    \end{tasks}

    \item 某元素失去两个电子后，在 $l=2$ 的轨道上刚好充满，则该元素为
    \begin{tasks}(4)
        \task \ce{Cr}
        \task \ce{Fe}
        \task \ce{Mn}
        \task \ce{Co}
    \end{tasks}

    \item 废弃的 \ce{CN-} 溶液不能倒入
    \begin{tasks}(2)
        \task 含 \ce{Fe^3+} 的废液中
        \task 含 \ce{Fe^2+} 的废液中
        \task 含 \ce{Cu^2+} 的酸性废液中
        \task 含 \ce{Cu^2+} 的碱性废液中
    \end{tasks}

    \item 下列固体中，与浓硫酸混合有 \ce{H2S} 生成的是
    \begin{tasks}(4)
        \task \ce{ZnF2}
        \task \ce{ZnCl2}
        \task \ce{ZnBr2}
        \task \ce{ZnI2}
    \end{tasks}

    \item 下列化合物中，在空气中可被氧化的是
    \begin{tasks}(4)
        \task \ce{Na2SO4}
        \task \ce{Mn(OH)2}
        \task \ce{Co2O3}
        \task \ce{MnO(OH)2}
    \end{tasks}

    \item 下列化合物中，碱性最强的是
    \begin{tasks}(4)
        \task \ce{N2H4}
        \task \ce{NH3}
        \task \ce{NH2OH}
        \task \ce{HN3}
    \end{tasks}

    \item 下列含氧酸中，属于一元酸的是
    \begin{tasks}(4)
        \task \ce{H3PO2}
        \task \ce{H3PO4}
        \task \ce{H3PO3}
        \task \ce{H2CO3}
    \end{tasks}

    \item 下列各组化合物中，颜色基本相同的一组是
    \begin{tasks}(2)
        \task \ce{Mn(OH)2}, \ce{MnO2}
        \task \ce{CoSO4 * 7H2O}, \ce{CoCl2 * 6H2O}
        \task \ce{CuSO4}, \ce{CuCl2}
        \task \ce{NiSO4 * 7H2O}, \ce{CoSO4 * 7H2O}
    \end{tasks}

    \item 下列水合晶体中，加热脱水时不生成碱式盐的是
    \begin{tasks}(4)
        \task \ce{MgCl2 * 6H2O}
        \task \ce{NiCl2 * 6H2O}
        \task \ce{CoCl2 * 6H2O}
        \task \ce{CuCl2 * 2H2O}
    \end{tasks}

    \item 下列配合物中，几何构型为正方形的是
    \begin{tasks}(4)
        \task \ce{[Cu(NH3)4]^2+}
        \task \ce{[Zn(NH3)4]^2+}
        \task \ce{[AuCl4]-}
        \task \ce{[Ni(CO)4]}
    \end{tasks}
\end{enumerate}

\vspace{0.5cm}
\noindent\textbf{二、简答题（第1小题15分，2\(\sim\)7小题每题5分，共45分）}

\begin{enumerate}[label=\arabic*., itemsep=1.5ex]
    \item 解释实验现象并写出相关的化学反应方程式
    \begin{enumerate}[label=(\arabic*), noitemsep, topsep=0pt, partopsep=0pt]
        \item 向粉红色 \ce{CoCl2} 溶液滴加浓盐酸则溶液变蓝，加水稀释后又变成粉红色。
        \item 向稀的 \ce{HCl} 和 \ce{Na2SO3} 混合溶液中通入 \ce{H2S}，溶液变浑浊；向稀的 \ce{HCl} 和 \ce{Na2SO4} 混合溶液中通入 \ce{H2S}，溶液无变化。
        \item 在煤气灯上加热 \ce{NaNO3} 晶体时，没有棕色气体生成；若 \ce{NaNO3} 晶体中混有 \ce{MgSO4} 晶体时，则有棕色气体生成。
    \end{enumerate}

    \item \ce{[CoCl4]^2-} 和 \ce{[NiCl4]^2-} 为四面体结构，而 \ce{[CuCl4]^2-} 和 \ce{[PtCl4]^2-} 却为平面四边形结构，试说明原因。

    \item 为什么 \ce{CCl4} 遇水不水解，而 \ce{SiCl4} 和 \ce{NCl3} 易水解？写出水解反应方程式。

    \item 用三种方法鉴别 \ce{MnO2} 和 \ce{CuO}。

    \item 1986年美国化学家克里斯特使用 \ce{KMnO4}、\ce{HF}、\ce{KF}、\ce{H2O2}、\ce{SbCl5} 为原料，用化学法制得单质 \ce{F2}，适用反应方程式表示各步反应过程。

    \item 设计方案：分离溶液中的 \ce{Al^3+}、\ce{Zn^2+}、\ce{Cr^3+}、\ce{Mn^2+}。

    \item 比较 \ce{HClO}、\ce{HClO3}、\ce{HClO4} 的酸性和氧化性强弱，说明原因。
\end{enumerate}

\vspace{0.5cm}
\noindent\textbf{三、推断题（10分）}

橙红色固体 A 受热后生成深绿色的固体 B 和无色的气体 C，加热时 C 能与镁反应生成固体 D。固体 B 与 \ce{NaOH} 共熔后溶于水生成绿色的溶液 E，在 E 中加入适量的 \ce{H2O2} 则生成黄色溶液 F，将 F 酸化变成橙色的溶液 G，在 G 中加 \ce{BaCl2} 溶液，得黄色沉淀。在 G 的浓溶液中加 \ce{KCl} 固体，溶解完全后，经蒸发、浓缩、冷却有橙红色晶体 H 析出，H 受强热得到的固体产物中有 B，同时得到能支持燃烧的气体 I，写出 A,B,D,E,H,I 所代表的物质的化学式，并写出由 E 到 F，H 到 I 过程的反应方程式。

\vspace{1cm}
\begin{center}
    {\Large \textbf{物理化学部分(共75分)}}
\end{center}

\vspace{0.5cm}
\noindent\textbf{一、选择题(24 pts.)}

\begin{enumerate}[label=\arabic*., itemsep=1.5ex]
    \item 用同一个电导池分别测定浓度为 $0.01\,\mathrm{mol/kg}$ 和 $0.1\,\mathrm{mol/kg}$ 的两个电解质溶液，其电阻分别为 $1000\,\Omega$ 和 $500\,\Omega$，则它们的摩尔电导率之比为( )
    \begin{tasks}(4)
        \task $1:5$
        \task $5:1$
        \task $10:5$
        \task $5:10$
    \end{tasks}

    \item 对于同一电解质的水溶液，当其浓度逐渐增加时，何种性质将随之增加( )
    \begin{tasks}(2)
        \task 在稀溶液范围内的电导率
        \task 摩尔电导率
        \task 电解质的离子平均活度系数
        \task 离子淌度
    \end{tasks}

    \item 某放射性同位素的半衰期为 $5\,\mathrm{d}$，则经 $15\,\mathrm{d}$ 后，所剩的同位素的量是原来的( )
    \begin{tasks}(4)
        \task $1/3$
        \task $1/4$
        \task $1/8$
        \task $1/16$
    \end{tasks}

    \item $\ce{AgCl}$ 在以下溶液中的溶解度递增次序为( )
    \begin{enumerate}
        \item $0.1\,\mathrm{mol/L}\ \ce{NaNO3}$
        \item $0.1\,\mathrm{mol/L}\ \ce{NaCl}$
        \item $\ce{H2O}$
        \item $0.1\,\mathrm{mol/L}\ \ce{Ca(NO3)2}$
        \item $0.1\,\mathrm{mol/L}\ \ce{NaBr}$
    \end{enumerate}
    \begin{tasks}(4)
        \task a$<$b$<$c$<$d$<$e
        \task b$<$c$<$a$<$d$<$e
        \task c$<$a$<$b$<$e$<$d
        \task c$<$b$<$a$<$d$<$e
    \end{tasks}

    \item 电导测定应用广泛，但下列问题中哪个是不能用电导测定来解决的？( )
    \begin{tasks}(2)
        \task 求难溶盐的溶解度
        \task 求弱电解质的解离度
        \task 求平均活度系数
        \task 测电解质溶液的浓度
    \end{tasks}

    \item 已知反应 $3\ce{O2(g)}=2\ce{O3(g)}$，在 $25^\circ\text{C}$ 时，$\Delta_r H_m^\ominus=-280\,\mathrm{J/mol}$，则对该反应有利的条件是( )
    \begin{tasks}(4)
        \task 升温升压
        \task 升温降压
        \task 降温升压
        \task 降温降压
    \end{tasks}

    \item 在 $1100^\circ\text{C}$ 时，发生下列反应
    \begin{enumerate}
        \item $\ce{C(s) + 2S(s) = CS2(g)}$ \quad $K_1=0.258$
        \item $\ce{Cu2S(s) + H2(g) = 2Cu(s) + H2S(g)}$ \quad $K_2=3.9\times10^{-3}$
        \item $2\,\ce{H2S(g) = 2H2(g) + 2S(s)}$ \quad $K_2=2.29\times10^{-2}$
    \end{enumerate}
    则 $1100^\circ\text{C}$ 时反应 $\ce{C(s) + 2Cu2S(s) = 4Cu(s) + CS2(g)}$ 的 $K$ 为
    \begin{tasks}(4)
        \task $8.99\times10^{-8}$
        \task $8.99\times10^{-5}$
        \task $3.69\times10^{-5}$
        \task $3.69\times10^{-8}$
    \end{tasks}

    \item 固体氧化物的分解压（分解反应是吸热的），当温度升高时( )
    \begin{tasks}(2)
        \task 分解压降低
        \task 分解压增大
        \task 分解压恒定
        \task 分解压不能确定
    \end{tasks}

    \item 同一个反应在相同反应条件下未加催化剂时平衡常数及活化能为 $K$ 及 $E_a$，加入正催化剂后则为 $K'$，$E_a'$，只存在下述关系( )
    \begin{tasks}(2)
        \task $K=K',\ E_a=E_a'$
        \task $K\neq K',\ E_a\neq E_a'$
        \task $K=K',\ E_a>E_a'$
        \task $K>K',\ E_a>E_a'$
    \end{tasks}

    \item 已知分解反应 $\ce{NH2COONH4(s)=2NH3(g) + CO2(g)}$，在 $50^\circ\text{C}$ 时的平衡常数 $K=6.55\times10^{-4}$，则此时 $\ce{NH2COONH4(s)}$ 的分解压力为( )
    \begin{tasks}(4)
        \task $16.63\times10^3\,\mathrm{Pa}$
        \task $5.940\times10^3\,\mathrm{Pa}$
        \task $5.542\times10^3\,\mathrm{Pa}$
        \task $2.928\times10^3\,\mathrm{Pa}$
    \end{tasks}

    \item 恒压条件下，加入惰性气体能增大哪一个反应的平衡转化率( )
    \begin{tasks}(1)
        \task $\ce{C6H5C2H5(g) = C6H5C2H3(g) + H2(g)}$
        \task $\ce{CO(g) + H2O(g) = CO2(g) + H2(g)}$
        \task $3/2\,\ce{H2(g) + 1/2N2(g) = NH3(g)}$
        \task $\ce{CH3COOH(l) + C2H5OH(l) = H2O(l) + C2H5COOCH3(l)}$
    \end{tasks}

    \item 若算得电池反应的电池电动势为负值时，表示此电池反应是( )
    \begin{tasks}(2)
        \task 正向进行
        \task 逆向进行
        \task 不可能进行
        \task 反应方向不能确定
    \end{tasks}
\end{enumerate}

\vspace{0.5cm}
\noindent\textbf{五、计算(10 pts.)}

取 $0^\circ\text{C}$、$3p^\ominus$ 的 $\ce{O2(g)}\ 10\,\mathrm{dm^3}$，绝热膨胀到压力为 $p^\ominus$，分别计算下列两过程的 $\Delta G$。
\begin{enumerate}
    \item 绝热可逆膨胀
    \item 将外压力骤减至 $p^\ominus$，气体反抗外压力进行绝热不可逆膨胀，假定 $\ce{O2}$ 为理想气体，其 $C_{V,\text{m}}=5/2R$，氧气的摩尔标准熵 $S_m^\ominus(298\,\mathrm{K})=205.0\,\mathrm{J/(mol\cdot K)}$
\end{enumerate}

\vspace{0.5cm}
\noindent\textbf{六、计算(15 pts.)}

氢醌的蒸汽压实验数据如下：

\begin{center}
\begin{tabular}{|c|cc|cc|}
\hline
 & \multicolumn{2}{c|}{固体 $\rightleftharpoons$ 气体} & \multicolumn{2}{c|}{液体 $\rightleftharpoons$ 气体} \\
\hline
$T/\mathrm{K}$ & 405.6 & 436.7 & 465.2 & 489.7 \\
$P/\mathrm{kPa}$ & 0.13334 & 1.3334 & 5.3327 & 13.334 \\
\hline
\end{tabular}
\end{center}

求：
\begin{enumerate}
    \item 氢醌的摩尔升华焓、摩尔蒸发焓、摩尔熔化焓
    \item 氢氧气、液、固三相平衡共存的温度和压力
\end{enumerate}

\vspace{0.5cm}
\noindent\textbf{七、计算(13 pts.)}

设 $\ce{N2O5(g)}$ 的分解为基元反应，在不同的温度下测得的速率常数 $k$ 值如下表所示：

\begin{center}
\begin{tabular}{|c|c|c|c|c|}
\hline
$T/\mathrm{K}$ & 273 & 298 & 318 & 338 \\
\hline
$k/\mathrm{min^{-1}}$ & $4.7\times10^{-5}$ & $2.0\times10^{-3}$ & $3.0\times10^{-2}$ & $0.30$ \\
\hline
\end{tabular}
\end{center}

试从这些数据求：阿仑尼乌斯经验式中的指数因子 $A$，实验活化能 $E_a$，$273\,\mathrm{K}$ 时过渡态理论中的 $\Delta^\ddagger H_\text{m}^\ominus$、$\Delta^\ddagger S_\text{m}^\ominus\ (c^\ominus)$

\vspace{0.5cm}
\noindent\textbf{八、计算题(13 pts.)}

电池 $\ce{Ag|AgCl(s)|Cl^-(a=0.8)|Cl2(P^*)|Pt}$ 在 $298\,\mathrm{K}$ 时的电动势 $E=1.1365\,\mathrm{V}$。
已知 $(\partial E/\partial T)_p=-5.993\times10^{-4}\,\mathrm{V/K}$

\begin{enumerate}
    \item 写出电极反应式和电池反应式。
    \item 计算该电池在 $298\,\mathrm{K}$ 时的 $\Delta_r G_m^\ominus$、$\Delta_r H_m^\ominus$、$\Delta_r S_m^\ominus$ 和标准平衡常数 $K$。
\end{enumerate}